\section*{Spivak Capítulo 3, Problema 20}
\addcontentsline{toc}{section}{Spivak Capítulo 3, Problema 20}

El ejercicio tiene dos apartados relacionados con condiciones de
continuidad/contracción.

\textbf{(a)} "Hallar una función $f$ que no sea constante y tal que
$|f(y)-f(x)|\le |y-x|$ para todos $x,y$." 

Una elección natural es cualquier función Lipschitz con constante 1; por
supuesto las funciones lineales $f(x)=cx$ con $|c|\le 1$ cumplen la
condición, y no son constantes cuando $c\neq0$. Por ejemplo, podemos tomar
\[ f(x)=x \qquad\text{o} \qquad f(x)=\sin x.\]
Ambas satisfacen $|f(y)-f(x)|\le |y-x|$; la primera lo hace por igualdad.
En el caso de $\sin x$ se puede justificar de varias maneras: por ejemplo,
el teorema del valor medio garantiza que existe algún $c$ entre $x$ y $y$ tal
que
\[\sin y-\sin x = (y-x)\cos c.\]
Como $|\cos c|\le1$ se deduce inmediatamente
$|\sin y-\sin x|\le |y-x|$. Alternativamente, empleando la identidad trigonométrica
$\sin y-\sin x = 2\cos\bigl(\tfrac{x+y}{2}\bigr)\sin\bigl(\tfrac{y-x}{2}\bigr)$
se obtiene
\[|\sin y-\sin x| = 2\bigl|\cos\tfrac{x+y}{2}\bigr|\,\bigl|\sin\tfrac{y-x}{2}\bigr|
\le 2\cdot1\cdot\Bigl|\tfrac{y-x}{2}\Bigr| = |y-x|,
\]
pues $|\sin t|\le|t|$ para todos los reales $t$. Cualquiera de estas
justificaciones pertenece al análisis elemental y muestra que la desigualdad
es válida para todas las parejas $x,y$.

También se pueden elegir funciones fractales o cono subyacente siempre que
la pendiente en valor absoluto sea como máximo 1 en cada punto; el punto
clave es que la desigualdad es exactamente la condición Lipschitz con
constante 1.

\textbf{(b)} Supón ahora que
$$f(y)-f(x)\le (y-x)^2\qquad\text{para todo }x,y.\tag{*}$$
Primero notemos que esta desigualdad forzaría también
$|f(y)-f(x)|\le (y-x)^2$ (la pregunta lo observa). Efectivamente, si escribimos
(*) con $x$ y $y$ intercambiados obtenemos
$$f(x)-f(y)\le (x-y)^2=(y-x)^2,$$
y combinando ambas desigualdades se obtiene la desigualdad en valor
absoluto mencionada.

Para probar que $f$ es constante partimos de (*) y aplicamos una técnica de
particionamiento sugerida en la indicación: dado $x<y$, dividamos el
intervalo $[x,y]$ en $n$ subintervalos iguales. Es decir, define
$$x_k = x + k\frac{y-x}{n}\qquad (k=0,1,\dots,n).$$
Entonces $x_0=x$ y $x_n=y$, y además cada diferencia $x_{k+1}-x_k=(y-x)/n$.
Aplicando (*) con $x=x_k$ y $y=x_{k+1}$ obtenemos
$$f(x_{k+1})-f(x_k) \le (x_{k+1}-x_k)^2 = \frac{(y-x)^2}{n^2}.	ag{1}$$
Sumando las desigualdades (1) para $k=0,\ldots,n-1$ se telescopa el lado
izquierdo y hallamos
$$f(y)-f(x) \le n\cdot\frac{(y-x)^2}{n^2} = \frac{(y-x)^2}{n}.	ag{2}$$
La clave es que la cota derecha puede hacerse arbitrariamente pequeña si
eligimos $n$ grande. Puesto que la desigualdad vale para todo $n$, enviamos
$n\to\infty$ y deducimos
$$f(y)-f(x) \le 0$$
en el caso $y>x$. Invirtiendo los papeles de $x$ e $y$ obtenemos
$$f(x)-f(y) \le 0$$
siempre. Por lo tanto $f(y)-f(x)=0$ para todos $x,y$, lo que implica que
$f$ es constante.

Esta es la demostración solicitada. En resumen, la condición cuadrática
implica una contracción tan fuerte que obliga a la función a ser rígida.

(No se necesita construir contrapartidas no constantes en la parte (b), la
conclusión única es la constancia.)